% Style file: https://media.neurips.cc/Conferences/NeurIPS2025/Styles.zip
\documentclass{article}
\usepackage[preprint]{neurips_2025}
\usepackage{hyperref}
\usepackage{url}
\usepackage{booktabs}
\usepackage{amsfonts}
\usepackage{amsmath}
\usepackage{nicefrac}
\usepackage{microtype}
\usepackage{graphicx}
\usepackage{natbib}
\usepackage{algorithm}
\usepackage{algorithmic}
\usepackage{adjustbox}
\usepackage[utf8]{inputenc}
\usepackage[T1]{fontenc}
\usepackage{lmodern}

\title{GOLF: A Negative Result for Codec-Aware QAT in Lane\\_3}

\author{Anonymous}

\begin{document}
\maketitle

\begin{abstract}
Deployment-constrained language modeling in Parameter Golf lane\\_3 requires optimizing validation bits-per-byte (val\\_bpb) after an int8-plus-zlib roundtrip under a 16 MB artifact cap and a $\leq$1.10$\times$ runtime gate. Prior work on quantization-aware training and model compression typically optimizes float-domain losses at fixed precision, leaving entropy coding as a post-hoc step that is not reflected in the training objective. This paper presents a single-condition pilot study of GOLF, a staged quantization-aware training setup designed for lane\\_3 that activates fake int8 quantization halfway through training and optionally adds a lightweight codec-inspired regularizer, while relying on the official Parameter Golf harness for the final int8-plus-zlib evaluation. Using the official float-trained baseline metrics as reference, we run one GOLF configuration and obtain a post-roundtrip val\\_bpb of 3.3919 compared to the baseline's 3.3794, and an evaluation runtime of 80894.3801 ms compared to 62992.0 ms, corresponding to a runtime ratio of 1.2842 and a failed quick-gate. This single-run negative result shows that naïve staged QAT under the current wiring and hyperparameters yields slightly worse compression and substantially slower evaluation than the provided baseline. The study highlights that codec-aware QAT in lane\\_3 is non-trivial and that careful attention to implementation wiring, proxy design, and gate-aware overheads is essential before more ambitious methods can claim benefits.
\end{abstract}

\section{Introduction}

\label{sec:introduction}

Language models intended for deployment on constrained hardware must satisfy multiple objectives simultaneously, balancing predictive quality with storage, runtime, and energy constraints. Mobile devices, edge accelerators, and serverless environments, for instance, often limit the size of model artifacts and impose tight latency budgets, making it insufficient to report only float-domain perplexity or loss. Instead, deployment-oriented evaluations focus on the behavior of the compressed artifacts that are actually shipped and executed, including quantization and entropy coding. The Parameter Golf challenge formalizes this perspective for language modeling by defining lane\\_3, in which models are evaluated on validation bits-per-byte (val\\_bpb) after an int8 quantization and zlib compression--decompression roundtrip, subject to a 16 MB artifact-size cap and a quick-gate that restricts runtime overhead to at most 10\% relative to a fixed baseline. This setting transforms model design into a resource-aware optimization problem in which training procedures, quantization schemes, and compression behavior jointly determine deployment viability.

Existing approaches to quantization-aware training (QAT) and model compression offer powerful tools for reducing precision and footprint, but they typically do not account for generic entropy coders such as zlib in a principled way. Standard QAT methods insert fake quantization operators into the forward pass and use straight-through estimators during backpropagation to help models tolerate integer arithmetic, often targeting int8 or lower precision for efficient deployment on accelerators. Model compression pipelines routinely combine pruning, quantization, and simple entropy coding, and learned image compression methods optimize rate--distortion objectives that explicitly trade off reconstruction error and estimated code length. However, these techniques usually employ differentiable entropy models tailored to arithmetic coding rather than to dictionary-based codecs like zlib, and in large language models much of the deployment literature emphasizes post-training quantization or mixed-precision arithmetic without explicitly modeling downstream compression. As a result, there remains a conceptual and empirical gap between training-time objectives and the actual deployed path of int8-plus-zlib in lane\\_3.

Parameter Golf lane\\_3 occupies an interesting middle ground between classical quantization benchmarks, which focus on accuracy at fixed precision, and systems benchmarks such as MLPerf Tiny, which emphasize end-to-end performance under hardware constraints \cite{banbury2021mlperf}. In lane\\_3, the primary metric, val\\_bpb, effectively measures how well a compressed model predicts bytes of text, while the artifact-size and runtime gates enforce hard resource limits. This design resembles performance-per-resource metrics proposed in other domains, such as medical imaging \cite{selvan2024pepr} and efficient training analyses \cite{desislavov2021compute}, but with a distinctive twist: the deployment path includes both fixed-function int8 quantization and a general-purpose zlib compressor whose behavior depends on the serialized patterns of model parameters. Aligning training with this path might in principle yield models whose weights are not only quantization-friendly but also more compressible by zlib, yet the interaction between these components is poorly understood. There is limited published evidence on whether codec-aware objectives can improve val\\_bpb under strict runtime gates in language models, or whether the standard float-trained baseline remains hard to beat.

This paper presents GOLF as an initial attempt to explore staged QAT under the constraints of Parameter Golf lane\\_3. Conceptually, GOLF introduces fake quantization modules into a transformer-like language model and activates them partway through training according to a normalized progress schedule, so that the model can adapt to int8-like perturbations. Configuration flags also allow a codec-inspired regularizer to be added to the loss, intended as a lightweight proxy for zlib-friendly weight patterns, though in the current run the exact proxy and wiring remain opaque at the logging level. Crucially, training is conducted via the existing \texttt{train\_gpt.py} entrypoint with environment variables specifying QAT-related knobs, while the official Parameter Golf harness performs a single post-hoc int8-plus-zlib roundtrip and measures val\\_bpb and runtime. This separation ensures compatibility with the challenge infrastructure but also means that ``codec alignment'' is currently limited to training-time proxies and one final roundtrip rather than full end-to-end differentiable feedback.

The contributions of this paper are deliberately scoped to match the evidence available from a single pilot run. First, it provides a deployment-aligned problem formulation for lane\\_3 that connects val\\_bpb, artifact size, and runtime constraints into a coherent objective, clarifying what it means to optimize for int8-plus-zlib deployment in this setting. Second, it describes a minimal staged QAT setup for lane\\_3 that can be enabled purely through configuration around the canonical training script, including the main QAT-related hyperparameters and their values in the evaluated run. Third, it reports a negative empirical result from one GOLF configuration, showing that under the current wiring and hyperparameters, staged QAT with optional codec-inspired regularization produces slightly worse val\\_bpb and substantially higher runtime than the official float-trained baseline, thereby violating the quick-gate. Finally, it distills methodological lessons from this outcome, emphasizing the importance of honest, descriptive reporting in n=1 experiments and highlighting how configuration wiring and proxy design can limit the practical benefits of codec-aware QAT in lane\\_3.

\section{Related Work}

\label{sec:related_work}

Research on efficient language model deployment spans quantization, compression, and benchmarking, yet few works explicitly target the combination of int8 quantization and generic codecs under strict runtime gates as in Parameter Golf lane\\_3. This section surveys relevant threads and contrasts them with the present single-run case study.

Work on quantization-aware training has established that inserting fake quantization operations into the forward pass can preserve model accuracy at low precision. Early approaches for convolutional networks and mobile architectures used per-tensor or per-channel affine quantization with straight-through estimators to approximate gradients through the quantizer. Subsequent methods improved stability and calibration through range learning, percentile clipping, and layer-wise scaling. In transformer-based language models, QAT has been shown to enable int8 and even lower-precision inference with limited degradation in perplexity when carefully tuned. Unlike these works, which primarily aim to maintain float-domain perplexity or accuracy at reduced precision, our study focuses on val\\_bpb after an int8-plus-zlib roundtrip and demonstrates that a naïve staged QAT configuration does not necessarily improve this deployment-specific metric in lane\\_3.

Post-training quantization (PTQ) offers an alternative to QAT by calibrating quantizers using a small dataset without retraining the entire model. Methods such as GPTQ and AWQ apply matrix-wise quantization with error compensation to large language models and have reported strong perplexity preservation at int8 or lower precision. In contrast to QAT, PTQ operates on a frozen float model and typically does not alter the training process. While these methods are highly relevant for deployment, they usually measure success in terms of perplexity or downstream accuracy rather than val\\_bpb after generic entropy coding, and they rarely integrate zlib or similar compressors into the evaluation. Our GOLF configuration, by contrast, modifies training through staged QAT and leaves the final int8-plus-zlib evaluation to the Parameter Golf harness, highlighting that even when QAT is enabled, the baseline float-trained model with post-hoc quantization and compression can remain competitive---or superior---on the lane\\_3 metric.

Learned compression literature explores rate--distortion trade-offs by integrating differentiable entropy models into the training objective. Balle et al. introduced end-to-end optimized image compression networks with continuous relaxations of quantization and hyperprior-based entropy models, allowing gradients to flow through approximate rate estimates. Minnen et al. improved these methods via context models that better capture local spatial dependencies in latent codes. While these approaches provide a template for codec-aware training, they typically rely on arithmetic coding and learned probability models rather than fixed-function compressors. In contrast, Parameter Golf lane\\_3 enforces int8 quantization followed by zlib, whose LZ77 and Huffman components exploit repeated substrings and symbol frequencies in serialized parameter bytes. Our single GOLF run attempts to approximate codec awareness via a lightweight regularizer, but unlike the learned compression works, it does not define a fully specified entropy model nor does it backpropagate through an accurate rate estimate.

Benchmarks such as MLPerf Tiny and other deployment-oriented suites emphasize end-to-end performance metrics that combine accuracy with latency or energy consumption \cite{banbury2021mlperf}. MLPerf Tiny, for example, defines microcontroller-focused tasks and mandates standardized hardware and reference implementations, reporting throughput and energy alongside accuracy. Similarly, studies on compute trends in deep learning advocate for metrics that reflect actual training and inference costs, not just parameter counts \cite{desislavov2021compute}. Parameter Golf lane\\_3 shares this deployment-centric spirit by requiring models to satisfy an artifact-size cap and a runtime quick-gate, and by defining val\\_bpb after the deployment roundtrip as the core metric. Unlike MLPerf-style work, which often evaluates mature, well-tuned systems, our paper reports an early negative result from a single staged QAT configuration, illustrating that even conceptually aligned methods may fail under strict gates if their overheads are not carefully controlled.

Compression and quantization for large language models have recently received considerable attention in the context of serving and inference cost reduction. Techniques such as weight-only quantization, activation quantization, and mixed-precision tensor formats are now widely used to accelerate inference on GPUs and specialized hardware. Some works combine these techniques with structured pruning or low-rank adaptation to further shrink models. These efforts, however, usually focus on latency and throughput rather than on the compressibility of serialized model artifacts under general-purpose codecs. In contrast, lane\\_3's combination of val\\_bpb and zlib compression foregrounds storage efficiency alongside predictive quality, and our GOLF run shows that simply adding QAT without carefully validated codec proxies can degrade both compression and runtime relative to the baseline.

Methodological work on evaluation and benchmarking in machine learning emphasizes the importance of strong baselines, transparent ablations, and honest reporting of negative results. Teterwak et al. showed that well-tuned empirical risk minimization baselines can match or outperform many domain generalization methods when optimization effort is controlled, arguing for rigorous comparisons. Liu et al. proposed the BOND benchmark for graph anomaly detection, highlighting how multi-metric evaluation can reveal trade-offs between detection performance and efficiency. Selvan et al. advocated performance-per-resource metrics in medical imaging, suggesting that resource-normalized scores better reflect deployment considerations \cite{selvan2024pepr}. Our paper follows this methodological direction by explicitly comparing a single GOLF configuration against the official float-trained baseline under identical lane\\_3 metrics and by reporting that the candidate is both slightly worse in val\\_bpb and substantially slower in runtime.

Several works in computer vision and remote sensing discuss benchmark design and deployment-aware evaluation but are less directly connected to codec-aware QAT. For example, AiTLAS provides a unified framework for training and evaluating models on earth observation data, while other benchmarks address robust translation, multimodal understanding, or manipulation detection. Unlike these benchmarks, which diversify tasks and modalities, Parameter Golf focuses narrowly on compressible language modeling under strict artifact and runtime gates. Our study draws on their general lesson that standardized evaluation harnesses are valuable, yet our actual implementation remains minimal: a single wrapper around the canonical training script with QAT-related environment variables, and reliance on the official harness for one post-hoc int8-plus-zlib evaluation.

Taken together, prior work on QAT, PTQ, learned compression, and deployment benchmarks provides a conceptual foundation for codec-aware training but does not directly address the precise combination of int8-plus-zlib and quick-gate constraints found in lane\\_3. In contrast to methods that report positive results across many ablations and seeds, our contribution is a focused negative case study: a single GOLF configuration that activates staged QAT via configuration in the Parameter Golf codebase and demonstrably fails to outperform the float-trained baseline on the lane\\_3 metrics. This outcome underscores that even well-motivated techniques from the broader literature require careful adaptation and validation to be effective in deployment-constrained settings.

\section{Method}

\label{sec:method}

The methodological goal of this work is to define and implement a minimal staged quantization-aware training setup for Parameter Golf lane\\_3, while preserving the canonical training entrypoint and using the official harness for final evaluation. This section first formalizes the lane\\_3 objective and then describes the specific QAT configuration and wiring used in the single GOLF run.

\subsection{Lane\\_3 objective and deployment operator}

\label{sec:lane_3_objective_and_deployment_operator}

Let \( f_\theta \) denote a transformer-style language model with parameters \( \theta \in \mathbb{R}^P \) trained on a corpus of token sequences. The float-domain training objective is the standard autoregressive negative log-likelihood over tokens, averaged over sequence positions. For deployment in lane\\_3, however, the model is transformed via an int8 quantization followed by zlib compression and decompression, and the evaluation metric is validation bits-per-byte (val\\_bpb) on text measured after this full roundtrip.

The deployment path can be represented as a composition of operators applied to \( \theta \). First, an int8 quantizer \( Q \) maps float parameters to signed 8-bit integers. Second, a serialization operator converts these integers into a contiguous byte sequence. Third, zlib compresses this sequence into a shorter bitstream and can later decompress it back to the original bytes. Finally, a dequantization operator reconstructs approximate float parameters from the decompressed integers. The combined deployment operator \( \mathcal{D} \) thus maps \( \theta \) to deployed parameters \( \tilde{\theta} = \mathcal{D}(\theta) \) that are used for inference in the evaluation harness.

Given a validation corpus with total byte length \(B\), val\\_bpb is computed by the Parameter Golf harness as the average negative log-likelihood of the deployed model normalized by \(B \log 2\), so that lower values correspond to better effective compression of the input text. The harness also records the total size in bytes of the compressed artifact produced by zlib and measures wall-clock runtime \(T(\theta)\) for a standardized validation workload. The lane\\_3 quick-gate compares the runtime ratio
\[
\mathrm{runtime\_ratio}(\theta) = \frac{T(\theta)}{T_\text{base}},
\]
where \(T_\text{base}\) is the baseline float-model runtime, and declares a run deployable only if \(\mathrm{runtime\_ratio}(\theta) \leq 1.10\). The artifact size must not exceed 16,000,000 bytes, but this check is handled internally by the challenge infrastructure.

In this work, we do not modify the deployment operator or the evaluation harness. Instead, we configure training to introduce QAT behavior in the base model while trusting the official Parameter Golf scripts to perform a single int8-plus-zlib evaluation at the end of training and to produce val\\_bpb, runtime, and quick-gate metrics. This separation ensures that any changes in lane\\_3 metrics arise from differences in training, not from altered evaluation logic.

\subsection{Staged quantization-aware training configuration}

\label{sec:staged_quantization_aware_training_confi}

The staged QAT behavior in GOLF is implemented through configuration knobs that are interpreted by the existing training code. The training script \texttt{train\_gpt.py} is used unmodified, and all QAT-related logic is activated via environment variables or configuration dictionaries passed into the canonical entrypoint. For the single GOLF run discussed here, the relevant hyperparameters are:

\begin{itemize}
 \item \texttt{qat\_enabled = 1.0}, indicating that QAT is turned on in the model construction.
 \item \texttt{fake\_quant\_strength = 1.0}, specifying that fake quantization is applied at full strength once active.
 \item \texttt{qat\_start\_frac = 0.5}, meaning that fake quantization is enabled only after 50\% of training progress.
 \item \texttt{qat\_group\_size = 32}, controlling the granularity at which weights are grouped for quantization.
 \item \texttt{quant\_error\_weight = 1.0}, influencing how quantization error is weighted in any internal loss proxy.
 \item \texttt{compression\_reg\_enabled = 1.0} and \texttt{compression\_reg\_lambda = 0.0001}, enabling a codec-inspired regularizer with a small coefficient.
 \item \texttt{seed = 1337.0}, fixing the random seed for this configuration.
\end{itemize}

When \texttt{qat\_enabled} is set, the model constructor inserts fake quantization modules on selected weight tensors, such as embeddings and output projections, so that forward passes can simulate the effects of int8 quantization. These modules expose a switch that depends on the normalized training progress \(t\), which is tracked by the training loop. With \texttt{qat\_start\_frac = 0.5}, the training run uses full-precision weights for the first half of optimization and then switches to applying fake quantization for the remaining steps. The \texttt{fake\_quant\_strength} parameter multiplies the contribution of quantized weights in this late-stage regime, and with a value of 1.0 it corresponds to fully adopting the fake-quantized parameters wherever QAT is enabled.

The \texttt{qat\_group\_size} parameter specifies how parameters are partitioned for quantization. With a value of 32, contiguous groups of 32 elements in the relevant dimension share quantization scale and zero-point, providing a compromise between per-tensor and per-element quantization. This group-wise scheme is expected to influence both numerical accuracy and the patterns seen by zlib in the serialized artifact, although the present study does not vary this setting.

The codec-inspired regularizer is controlled by \texttt{compression\_reg\_enabled} and \texttt{compression\_reg\_lambda}. When enabled, the training loop adds a penalty term of magnitude proportional to \texttt{compression\_reg\_lambda} that is intended to favor weight configurations that compress better under zlib. The concrete form of this proxy is not exposed in the logs for the run analyzed here, and no additional ablations are available that isolate its effect. As a result, the negative outcome reported in the Results section should be interpreted primarily as a consequence of staged QAT under this configuration, with the understanding that any effect of the codec-inspired regularizer is entangled with the QAT behavior.

A critical practical detail is that training is executed with \texttt{SKIP\_FINAL\_INT8\_ROUNDTRIP\_EVAL=1} set in the environment. This flag instructs the training script to skip any internal int8-plus-zlib evaluation it might otherwise perform at the end of training. Instead, the Parameter Golf harness separately loads the resulting checkpoint, applies the official quantization and zlib pipeline once, and measures val\\_bpb and runtime. Consequently, while the training side is ``quantization-aware'' in that fake quantization is active for part of optimization, the only actual int8-plus-zlib roundtrip used for decision making is the single post-training one executed by the harness.

\subsection{Integration with the Parameter Golf harness}

\label{sec:integration_with_the_parameter_golf_harn}

The GOLF configuration is integrated into the Parameter Golf infrastructure by registering it as a new condition in the lane\\_3 experiment registry. The corresponding script sets environment variables for QAT-related hyperparameters, including \texttt{QAT\_ENABLED}, \texttt{QAT\_START\_FRAC}, \texttt{QAT\_GROUP\_SIZE}, \texttt{FAKE\_QUANT\_STRENGTH}, \texttt{COMPRESSION\_REG\_ENABLED}, and \texttt{COMPRESSION\_REG\_LAMBDA}, and then invokes the canonical \texttt{train\_gpt.py} entrypoint. Training proceeds to completion, producing a checkpoint whose path is recorded in the lane\\_3 records folder.

After training, the official lane\\_3 evaluation harness is run exactly once for this configuration. The harness performs the int8 quantization and zlib compression--decompression roundtrip on the checkpoint, measures the compressed artifact size to enforce the 16 MB limit, executes the validation workload with the deployed parameters to compute val\\_bpb, and times the evaluation to obtain \texttt{candidate\_runtime\_ms}. It also compares this runtime to the precomputed baseline runtime \texttt{baseline\_runtime\_ms} and sets \texttt{quick\_gate\_passed} to 0.0 because the ratio is greater than 1.10. All these values are stored in JSON files in the records folder, from which the numbers in the Results section are drawn.

The baseline metrics, \texttt{baseline\_val\_bpb = 3.3794} and \texttt{baseline\_runtime\_ms = 62992.0}, are treated as fixed reference values supplied by the challenge infrastructure. The GOLF configuration's metrics, \texttt{val\_bpb = 3.3919} and \texttt{candidate\_runtime\_ms = 80894.3801}, are measured once by the harness. The difference in val\\_bpb, \texttt{delta\_bpb = 0.0125}, and the runtime ratio, \texttt{runtime\_ratio = 1.2842}, are then computed by the harness and logged.

Within this workflow, GOLF does not introduce any new evaluation scripts, ablation engines, or detailed runtime decomposition. The ``ablation'' functionality mentioned in earlier drafts corresponds to simple comparisons between hyperparameter dictionaries and metric outputs in the records folder and does not participate in the execution of this particular run. As a result, the empirical analysis in this paper rests entirely on a single training run with QAT enabled and the official baseline metrics, rather than on a suite of systematically varied ablations.

\section{Experiments}

\label{sec:experiments}

The experimental evaluation consists of comparing the official float-trained baseline metrics with those from a single GOLF configuration in Parameter Golf lane\\_3. This section outlines the task and metrics, summarizes the configuration of the two conditions, and states the hardware and protocol used to obtain the reported numbers.

\subsection{Task and evaluation metrics}

\label{sec:task_and_evaluation_metrics}

Lane\\_3 uses a next-token prediction task on a fixed language modeling corpus defined by the challenge organizers. Models are trained on a training split and evaluated on a validation split, both of which are tokenized using a prescribed tokenizer. The underlying dataset paths and exact tokenization parameters are encoded in the Parameter Golf repository and are not modified in this work.

The primary evaluation metric is validation bits-per-byte, val\\_bpb, computed after the int8-plus-zlib deployment roundtrip. The harness loads the deployed model, runs it over the validation set to compute token-level negative log-likelihoods, sums these losses, and divides by the total number of input bytes times \(\log 2\), yielding an effective bits-per-byte measure. Lower val\\_bpb indicates better compression performance at the byte level, which serves as the objective for lane\\_3 under the constraint that the compressed artifact does not exceed 16 MB.

The harness also measures evaluation runtime for the deployed model by timing the full validation workload. This runtime is reported as \texttt{candidate\_runtime\_ms} for the GOLF configuration and as \texttt{baseline\_runtime\_ms} for the baseline. The quick-gate condition is applied by computing the ratio \texttt{runtime\_ratio = candidate\_runtime\_ms / baseline\_runtime\_ms} and marking the run as failed if this exceeds 1.10. The harness additionally computes \texttt{delta\_bpb = val\_bpb - baseline\_val\_bpb}, which is positive when the candidate is worse than the baseline in val\\_bpb.

\subsection{Baseline and GOLF configurations}

\label{sec:baseline_and_golf_configurations}

The baseline condition corresponds to the official float-trained model provided in the lane\\_3 harness. Its val\\_bpb and runtime were measured previously by the challenge infrastructure and are supplied as constants in the current records: \texttt{baseline\_val\_bpb = 3.3794} and \texttt{baseline\_runtime\_ms = 62992.0}. This baseline uses post-hoc int8 quantization and zlib compression for evaluation, but its training was conducted without the QAT configuration described in this paper.

The GOLF condition uses the same underlying model architecture, dataset, and evaluation path as the baseline, but modifies training by enabling staged QAT through the hyperparameters listed in the Method section. In particular, QAT is turned on, fake quantization is activated after half of training progress, group size is set to 32, and a codec-inspired regularizer is enabled with a small coefficient. The random seed \texttt{seed = 1337.0} is used to initialize this run. No other aspects of the training script, optimizer, or data loader are changed.

After training the GOLF configuration once, the Parameter Golf harness performs an int8-plus-zlib evaluation and records the following metrics:

\begin{itemize}

\end{itemize}

These values are used verbatim in the Results section and are not adjusted, rounded, or aggregated.

\subsection{Hardware and protocol}

\label{sec:hardware_and_protocol}

Training and evaluation were executed on a single NVIDIA GeForce RTX 5070 Ti GPU with approximately 16 GB of VRAM. The operating system, CPU, and storage configuration follow the default environment of the Parameter Golf leaderboard instance and are not modified for this study. The training script \texttt{train\_gpt.py} is invoked once for the GOLF configuration, with QAT-related environment variables set as described earlier. The script runs to completion within the standard per-run time budget enforced by the challenge.

Once training finishes, the lane\\_3 evaluation harness is invoked for this configuration. The harness loads the checkpoint, applies int8 quantization and zlib compression--decompression to obtain deployed parameters, measures the compressed artifact size against the 16 MB cap, and runs the validation workload while timing its duration. The harness writes all metrics to JSON files in the records folder, including the baseline and candidate values used in the analysis.

No additional seeds, repetitions, or ablations were run for this study. As a result, the experiment consists of exactly two data points: the fixed baseline metrics and the single GOLF run's metrics.

\section{Results}

\label{sec:results}

The results compare the official float-trained baseline with the single GOLF configuration in terms of post-roundtrip val\\_bpb and evaluation runtime. Because there is exactly one recorded run for the GOLF condition and the baseline metrics are supplied as constants, the analysis is purely descriptive and does not involve aggregation or statistical inference.

Table 1 summarizes the metrics for the two conditions using the values reported by the Parameter Golf harness.

\begin{table}[h]
\centering
\caption{Lane\_3 metrics for the official float-trained baseline and the single staged-QAT GOLF configuration. All numbers are taken directly from the Parameter Golf records. Lower val\_bpb and runtime are better; the quick-gate requires runtime\_ratio PROT0 1.10.}
\begin{tabular}\{lcccc\}
\hline
Condition \& val\\_bpb \& candidate\\_runtime\\_ms \& runtime\\_ratio \& quick\\_gate\\_passed \\
\hline
Baseline \& 3.3794 \& 62992.0 \& 1.0 \& 1.0 \\
GOLF (staged QAT) \& 3.3919 \& 80894.3801 \& 1.2842 \& 0.0 \\
\hline
\end{tabular}
\end{table}

The baseline achieves a val\\_bpb of 3.3794 and an evaluation runtime of 62992.0 ms. The GOLF configuration yields a val\\_bpb of 3.3919 and a runtime of 80894.3801 ms. The difference in val\\_bpb, \texttt{delta\_bpb = 0.0125}, indicates that the GOLF model is slightly worse in terms of effective compression on the validation corpus than the float-trained baseline when both are evaluated after the int8-plus-zlib roundtrip. The runtime ratio of 1.2842 shows that the GOLF evaluation is approximately 28\% slower than the baseline, leading the harness to mark \texttt{quick\_gate\_passed = 0.0} for this run.

Figure 1 visualizes the val\\_bpb comparison between the two conditions. The baseline bar is lower, reflecting better compression, while the GOLF bar is higher, consistent with its positive delta in val\\_bpb.

\begin{figure}[t]
\centering
\includegraphics[width=0.95\columnwidth]{charts/fig_main_results_comparison.png}
\caption{Val\\_bpb comparison}
\label{fig:val_bpb_comparison}
\end{figure}
\textbf{Figure 1.} Validation bits-per-byte (val\\_bpb) after the int8-plus-zlib roundtrip for the official float-trained baseline and the single staged-QAT GOLF configuration. The baseline achieves val\\_bpb = 3.3794, while GOLF yields val\\_bpb = 3.3919, indicating a modest degradation in compression performance for the candidate.

Figure 2 depicts the runtime and runtime ratio, with a dashed line at 1.10 to indicate the quick-gate threshold. The baseline lies exactly at ratio 1.0 by definition, and the GOLF configuration exceeds the gate with a ratio of 1.2842.

\begin{figure}[t]
\centering
\includegraphics[width=0.95\columnwidth]{charts/fig_runtime_overhead_quick_gate.png}
\caption{Runtime comparison}
\label{fig:runtime_comparison}
\end{figure}
\textbf{Figure 2.} Evaluation runtime for the official baseline and staged-QAT GOLF configuration in lane\\_3. The baseline runtime is 62992.0 ms, while GOLF takes 80894.3801 ms, producing a runtime ratio of 1.2842. The horizontal dashed line marks the 1.10 quick-gate threshold, which the GOLF configuration fails.

A simple joint view of val\\_bpb versus runtime reinforces the same conclusion: the baseline occupies the more favorable region with lower val\\_bpb and faster evaluation, whereas the GOLF configuration is both slightly worse in val\\_bpb and substantially slower.

\begin{figure}[t]
\centering
\includegraphics[width=0.95\columnwidth]{charts/fig_bitrate_vs_runtime_tradeoff.png}
\caption{Bitrate--runtime trade-off}
\label{fig:bitrate_runtime_trade_off}
\end{figure}
\textbf{Figure 3.} Joint plot of val\\_bpb versus evaluation runtime for the baseline and staged-QAT GOLF configuration. The baseline lies closer to the lower-left corner, indicating better compression and faster runtime. The GOLF configuration lies up and to the right, showing that this particular staged-QAT setup does not improve the bitrate--runtime trade-off.

Given the magnitude of the runtime difference---around 18 seconds on a baseline of roughly 63 seconds---it is unlikely that timer noise alone accounts for the observed overhead on a single machine. However, without repeated measurements or multiple seeds, the experiment cannot quantify variability or support statistical claims beyond these descriptive comparisons. Within this purely descriptive framework, the empirical conclusion is straightforward: under the single staged-QAT configuration evaluated here, GOLF does not surpass the float-trained baseline in lane\\_3 and instead produces a slightly worse val\\_bpb and a runtime that violates the quick-gate.

\section{Discussion}

\label{sec:discussion}

The single-run comparison between the official baseline and the staged-QAT GOLF configuration provides a concrete example of how conceptually appealing methods can underperform in deployment-constrained settings when evaluated honestly. The intention behind enabling QAT in lane\\_3 is to align training with the int8 deployment regime, thereby reducing the mismatch between float-domain optimization and post-hoc quantization. Yet the recorded metrics show that the evaluated GOLF configuration yields a small degradation in val\\_bpb and a substantial increase in runtime relative to the baseline, despite sharing the same architecture and dataset.

One likely factor in this outcome is the overhead introduced by fake quantization modules and any associated regularization. Even though QAT is only activated after half of training progress, its presence in the second half of training and any additional loss computations can increase the per-step cost of optimization or the cost of the forward pass during evaluation. In a setting like lane\\_3, where the quick-gate enforces a tight runtime budget, these overheads can directly render a method non-deployable, regardless of its potential to reduce val\\_bpb. This trade-off echoes observations from other resource-aware benchmarks, where methods that reduce parameter counts or precision sometimes impose enough computational overhead to negate their practical benefits \cite{banbury2021mlperf, desislavov2021compute}.

Another contributing factor is the incomplete specification and validation of the codec-inspired regularizer. While the configuration enables this regularizer with a small coefficient, the lack of a documented proxy function and the absence of ablations that isolate its effect make it difficult to assess how much it contributes to the observed degradation in val\\_bpb. In contrast, learned compression methods in vision explicitly design and validate entropy models against measured bitrates. Without such validation, a heuristic proxy could encourage weight distributions that negatively interact with zlib's LZ77 and Huffman stages, offsetting any gains from improved quantization robustness.

The experiment also illustrates the importance of keeping claims proportional to evidence in deployment-focused studies. It can be tempting to describe an initial QAT configuration and a minimal wrapper around a training script as a full `\texttt{framework'' or }`blueprint,'' especially when the conceptual story aligns with contemporary trends in efficient LLM deployment. However, the lane\\_3 results demonstrate that even a seemingly mild change such as enabling staged QAT must be assessed directly against baseline metrics using the official harness, and that a single negative run already provides meaningful information: namely, that this particular configuration of QAT and codec proxies does not push the bitrate--runtime frontier forward in lane\\_3.

For practitioners considering QAT or codec-aware training in similar settings, the main lesson is practical rather than theoretical. Before investing extensively in elaborate proxy losses or complex training schedules, it is valuable to establish a simple, reproducible baseline and then introduce changes incrementally, using the same evaluation harness and metrics. Comparing each candidate directly to the baseline's val\\_bpb and runtime, as done here, can quickly reveal whether a configuration is moving in a promising direction or, as in this case, degrading both key metrics. This stepwise approach parallels recommendations in benchmarking and domain generalization work, where strong empirical baselines often serve as a more reliable reference than complex, lightly tested methods \cite{teterwak2023improved, liu2022bond}.

The limitations of the present study are also instructive. With n=1 for the GOLF run and fixed baseline metrics, there is no basis for statistical inference or robustness claims. Nonetheless, the magnitude of the runtime overhead and the consistent direction of the val\\_bpb difference provide enough signal to conclude that this specific configuration is not competitive under lane\\_3's rules. Future work can build on this descriptive starting point by running additional seeds, varying QAT start fractions and group sizes, and turning the codec-inspired regularizer on and off to better understand which components are harmful or helpful. Such work would move the story from a single negative case to a more systematic exploration of codec-aware QAT under strict deployment gates.

\section{Limitations}

\label{sec:limitations}

This study has several limitations that constrain the scope of its conclusions. First, the experimental evidence is based on exactly one GOLF training run and the fixed baseline metrics supplied by the Parameter Golf harness, so the analysis is purely descriptive and does not estimate variability across seeds or runs. Second, the codec-inspired regularizer is enabled in the configuration but its concrete form, evaluation schedule, and effect on training are not documented in the logs or isolated in ablations, making it unclear how much it contributes to the observed degradation in val\\_bpb. Third, the staged QAT configuration itself is only one point in a large hyperparameter space, defined by a single QAT start fraction, group size, and fake quantization strength; different settings might yield different trade-offs that are not explored here. Fourth, the underlying model architecture, dataset, and tokenizer are taken as fixed from the lane\\_3 environment, so the results may not generalize to other model sizes, text domains, or tokenization schemes that interact differently with int8 quantization and zlib compression. Finally, runtime measurements are taken from a single evaluation per condition on one hardware setup, and while the large observed difference suggests a genuine overhead, the study does not quantify runtime variability or examine how it might change across machines or software stacks.

\section{Conclusion}

\label{sec:conclusion}

This paper examined a single staged quantization-aware training configuration, GOLF, in the context of the Parameter Golf lane\\_3 challenge, where models are evaluated on validation bits-per-byte after an int8-plus-zlib roundtrip under strict artifact-size and runtime constraints. By enabling QAT through configuration around the canonical training script and relying on the official harness for evaluation, we obtained one candidate run whose metrics can be compared directly to the provided float-trained baseline. The results show that this staged-QAT configuration yields a val\\_bpb of 3.3919 versus the baseline's 3.3794 and an evaluation runtime of 80894.3801 ms versus 62992.0 ms, corresponding to a runtime ratio of 1.2842 and a failed quick-gate. In other words, under the wiring and hyperparameters used here, staged QAT did not improve the deployment-aligned metric and introduced substantial runtime overhead.

These findings highlight that, in deployment-constrained settings like lane\\_3, enabling QAT or codec-inspired regularization is not automatically beneficial and must be evaluated directly against the specific gates and metrics that define success. The fact that a carefully engineered baseline remains superior on both val\\_bpb and runtime emphasizes the value of strong reference models and honest reporting of negative results. From a methodological perspective, the study suggests that future work on codec-aware training for language models should prioritize three directions: first, clarifying and validating any codec proxies against actual zlib behavior; second, systematically exploring QAT schedules and group sizes with ablations that isolate their effects; and third, incorporating multiple seeds or repeated measurements to quantify variability in both val\\_bpb and runtime. By building on this initial negative case with more extensive experimentation and sharper proxies, subsequent research can more reliably determine whether and when codec-aware QAT can outperform strong float-trained baselines under the demanding constraints of Parameter Golf lane\\_3.

\bibliographystyle{plainnat}
\bibliography{references}

\end{document}
